\chapter{Introduction}
\label{ch:introduction}

\section{Background}
\label{sec:background}

Warehouses sit at the point where a promise made to a customer becomes a physical operation.
An order that has been accepted must be picked from storage, packed for transport and released
to a carrier, and it must complete that journey inside whatever delivery window was promised
when the order was taken. The growth of online retail has made this harder in a specific way:
order volumes have become larger and more volatile, individual orders have become smaller and
more numerous, and the delivery promises attached to them have become both shorter and more
differentiated \citep{boysen2019}. A warehouse today rarely serves a single homogeneous demand
stream. It serves several, each with its own urgency, and it must serve them from one shared
pool of people and equipment.

Of the activities carried out inside a warehouse, order picking has long been identified as
the most labour-intensive and the most expensive, frequently accounting for the majority of
total warehouse operating expenditure \citep{dekoster2007}. That observation has motivated a
large body of research into how picking should be organised: how items should be stored, how
pickers should be routed, how orders should be batched and how zones should be allocated
\citep{gu2007,vangils2018}. The dominant concern in that literature is the efficiency of the
picking process itself, measured in travel distance, throughput or handling time.

This thesis is concerned with a decision that sits one level above those questions and that
receives comparatively less attention: given a month's expected demand, \emph{how many people
should be assigned to each stage of the fulfilment process}, and \emph{in what order should the
work be released to them}. These are the two levers that an operations manager actually
controls on a monthly planning cycle. Layout and storage policy are fixed over that horizon;
headcount and operating discipline are not.

\section{Motivation}
\label{sec:motivation}

The subject of this thesis was motivated by discussions held during a summer internship at
Baobab Soluciones, a company working on optimisation and decision-support applications for
industrial clients. What became apparent in those conversations was a recurring gap between
the sophistication of the analytical tools available and the form in which operational
decisions are actually taken. A manager asked to commit to a headcount for the coming month
does not need a theoretically optimal solution to an idealised problem. They need a defensible
answer to a small number of concrete questions: how many people, at which stage, at what cost,
with what risk of missing the service commitments, and what would change if one more person
were hired. The research reported here was subsequently developed as an independent academic
project, and is described in these terms in the Conflict of Interest statement.

Two features of the problem make it resistant to the obvious analytical shortcut. The first is
that a warehouse is a queueing system, not a flow of independent tasks. Dividing total expected
workload by capacity per worker produces a ratio, and that ratio ignores the fact that work
arrives unevenly, that variability in handling time causes queues to form, and that queues
propagate downstream. A stage loaded to what appears on paper to be a comfortable eighty-five
per cent of capacity can still accumulate a backlog that no aggregate calculation predicts
\citep{hopp2011}.

The second is that service commitments are differentiated. When a warehouse promises one
delivery window to urgent orders and a much longer one to normal orders, an aggregate
service-level figure becomes actively misleading. A configuration can report ninety per cent
overall on-time performance while failing the urgent class catastrophically, simply because
the urgent class is a small share of volume. Any planning instrument that reports a single
service number can therefore recommend a configuration that would be unacceptable in practice.

Together these two features suggest that the monthly capacity decision and the sequencing
decision cannot sensibly be separated. How many people are needed depends on how the work is
ordered, because a sequencing rule that protects the urgent class changes the number of people
required to satisfy the urgent commitment. Conversely, whether a sequencing rule is worth
adopting depends on how tightly capacity is constrained, because a rule that makes a decisive
difference under pressure may make none at all when capacity is ample.

A third consideration motivated the shape of this study as much as the first two. Planning
decisions of this kind are not taken by the analyst who builds the model; they are taken by a
manager who must commit a budget and defend it. That has consequences for what counts as a
useful output. A recommendation of a specific headcount, unaccompanied by any account of why
that number rather than another, is difficult to act on and easy to dismiss --- particularly
when it differs from what a simpler calculation suggested. An instrument intended for this
decision therefore has to expose its reasoning: which stage is constraining the operation, which
service commitment is at risk, what the marginal worker would cost and what it would buy. This
requirement shaped the framework as much as the modelling considerations did, and it is why the
bottleneck diagnostic and the explicit marginal-capacity test are treated in this thesis as
components of the method rather than as reporting conveniences.

\section{Research problem}
\label{sec:research-problem}

The research problem addressed by this thesis is therefore:

\begin{quote}
\emph{How can discrete-event simulation support monthly warehouse workforce-capacity planning
and order sequencing under heterogeneous demand and differentiated service-level requirements,
in a way that is both quantitatively grounded and explainable to the manager who must act on
the recommendation?}
\end{quote}

Three aspects of this formulation deserve emphasis. \textbf{Monthly} places the problem at a
tactical planning horizon: the unit of decision is a full-time equivalent of labour capacity
for a month, not a shift roster or a daily assignment. \textbf{Heterogeneous} recognises that
orders differ not only in size but in the workload they impose, and that the same order can be
demanding at one stage and trivial at another. \textbf{Differentiated} makes the service
requirement class-specific rather than aggregate. \textbf{Explainable} is a requirement in its
own right, not a presentational afterthought: a recommendation that cannot be justified to the
person accountable for the headcount will not be used.

\section{Objectives}
\label{sec:objectives}

The general objective of this thesis is to design, implement and evaluate a simulation-based
decision-support framework that recommends a monthly warehouse workforce configuration and an
order sequencing policy under differentiated service-level and economic constraints.

This general objective is pursued through five specific objectives:

\begin{enumerate}
  \item To formalise the monthly capacity-planning problem for a three-stage fulfilment
        operation in which orders impose stage-differentiated workload.
  \item To implement an analytical capacity anchor and a bounded search over workforce
        candidates, evaluated by discrete-event simulation under common random numbers.
  \item To compare three order-sequencing policies against class-specific service floors and
        an explicit economic objective.
  \item To provide stage-level bottleneck diagnostics and an adaptive procedure that tests
        whether additional capacity at the constraining stage is economically justified.
  \item To demonstrate that a single framework can serve both retrospective and
        forecast-based planning, and to characterise how the two differ.
\end{enumerate}

\section{Research questions}
\label{sec:research-questions}

Four research questions operationalise these objectives. Each is answerable from the
experimental evidence generated by the implemented framework, and each is answered explicitly
in \Cref{ch:conclusions}.

\begin{description}
  \item[RQ1] How can discrete-event simulation support monthly workforce-capacity planning
    under heterogeneous warehouse demand, and what does it add over an analytical capacity
    estimate alone?
  \item[RQ2] How do alternative order-sequencing policies affect differentiated service-level
    attainment and operating cost, both at equal workforce and at each policy's own feasible
    workforce?
  \item[RQ3] How can stage-level bottleneck diagnostics and adaptive capacity evaluation
    improve the explainability and economic interpretation of workforce recommendations?
  \item[RQ4] How can the same simulation framework support both retrospective and
    forecast-based planning, and how do their uncertainty structures differ?
\end{description}

\Cref{tab:rq-map} maps each question to the method used to answer it and to the evidence
presented in \Cref{ch:results}.

\begin{table}[htbp]
\centering
\caption{Research questions, methods and evidence.}
\label{tab:rq-map}
\small
\begin{tabularx}{\textwidth}{lXXl}
\toprule
\textbf{RQ} & \textbf{Method} & \textbf{Principal evidence} & \textbf{Section} \\
\midrule
RQ1 & Analytical capacity estimate compared against simulated evaluation of 16 workforce
      candidates per month & Analytical estimate versus recommended workforce, all 12 months
      & \ref{sec:res-capacity} \\
RQ2 & Three policies evaluated on identical stochastic realisations across 576 month
      \(\times\) workforce \(\times\) policy runs & Feasibility counts, class-level service
      trade-off, representative capacity regimes & \ref{sec:res-policy} \\
RQ3 & Stage Pressure decomposition and the adaptive capacity search trail & Bottleneck ranking
      per month; accept/reject outcome of adding one FTE & \ref{sec:res-bottleneck} \\
RQ4 & Comparison of the retrospective and forecast-based workflows as executed & Replicated
      December planning run with feasibility probabilities & \ref{sec:res-future} \\
\bottomrule
\end{tabularx}
\end{table}

\section{Scope and delimitations}
\label{sec:scope}

The scope of this study is deliberately bounded, and the boundaries matter for interpreting
the results.

The modelled operation consists of three sequential stages --- picking, packing and dispatch
--- each with its own queue and its own pool of workers. Stages are not physically detailed:
there is no travel time, no aisle congestion, no storage layout and no finite buffer between
stages. The model represents the \emph{capacity and queueing} behaviour of the fulfilment
process, not its physical geography. Questions of layout, routing and storage assignment, which
form the bulk of the order-picking literature \citep{dekoster2007,weidinger2019}, are outside
this scope and are treated as fixed.

The planning unit is a monthly full-time equivalent. A recommendation of twenty-two picking
FTE means twenty-two months of full-time picking capacity, not twenty-two people standing on
the floor at the same moment. Shift design, rostering, breaks, absence and overtime are not
modelled. Converting an FTE recommendation into a roster is a separate problem that a planner
would address afterwards.

Time in the simulation is \emph{operating time}, not calendar time. A month of demand is mapped
onto the productive hours that the workforce is actually paid for, preserving the relative
ordering and burstiness of arrivals but not literal wall-clock timestamps. Service-level
thresholds are consequently measured in productive elapsed minutes. This is explained in detail
in \Cref{sec:operating-time}, and its consequences are revisited among the limitations.

Finally, all experiments reported here use synthetic data constructed to be structurally
realistic rather than data drawn from a real warehouse. The framework accepts real order-level
histories as input --- that is one of its two operating modes --- but the experimental evidence
in this thesis was generated from a synthetic dataset. \Cref{sec:data-provenance} states this
precisely, and the distinction between what the system \emph{can} consume and what these
experiments \emph{did} consume is maintained throughout.

\section{Contributions}
\label{sec:contributions}

This thesis makes three kinds of contribution.

\textbf{Methodologically}, it demonstrates a design in which the physical capacity available in
simulation and the capacity paid for in the economic model are forced to be the same quantity,
removing a discrepancy that would otherwise grant workers unpaid processing time; a
finite-horizon monthly formulation in which work left unfinished is retained as explicit
backlog rather than absorbed by unlimited post-period processing; a feasibility criterion that
is class-specific rather than aggregate; and a three-way policy comparison in which a learned
policy is evaluated on exactly the same stochastic realisations as two transparent heuristics.

\textbf{Practically}, it delivers a workflow that converts either an order-level history or an
aggregate forecast into a recommendation a manager can act on and interrogate: workforce per
stage, sequencing policy, expected cost, class-level service attainment, a named location of
operational pressure with its decomposition, and an explicit economic test of whether one more
full-time equivalent would pay for itself.

\textbf{Empirically}, it provides evidence that the value of sequencing policy is
\emph{capacity-regime dependent}. Under ample capacity the three policies studied are
practically indistinguishable; under pressure the choice determines whether the service
commitments can be met at all. In the four highest-demand months of the year studied, the
first-in-first-out discipline satisfied the service floors at none of the sixteen workforce
configurations examined, while both alternatives satisfied them at almost all. This reframes
the question from ``which sequencing policy is best?'' to ``under which conditions does
sequencing policy matter?''.

The thesis does \emph{not} claim to contribute a new reinforcement learning algorithm, a
globally optimal capacity solution, or a model validated against real industrial data. The
learning agent evaluated here is a standard deep Q-network \citep{mnih2015}; the search over
workforce configurations is a bounded structured neighbourhood, not an optimisation; and the
validation performed is internal rather than external.

\section{Structure of the thesis}
\label{sec:structure}

\Cref{ch:theory} reviews the theoretical foundations and the relevant literature, covering
warehouse fulfilment operations, capacity and workforce planning, queueing and congestion,
discrete-event simulation, sequencing and priority rules, differentiated service, and
reinforcement learning for operational control, and closes by positioning the contribution of
this work.

\Cref{ch:framework} formalises the problem and presents the decision-support framework: the
system under study, the notation used throughout, the architecture of the framework and the
two workflows it supports.

\Cref{ch:methodology} describes the methodology in detail --- the data and its provenance, the
operating-time abstraction, the simulation engine, the three sequencing policies as they are
implemented, the reinforcement learning specification, the capacity estimate and candidate
generation, the feasibility and economic criteria, the bottleneck diagnostic, the adaptive
capacity search, the experimental design including common random numbers, and the verification
and validation performed.

\Cref{ch:results} reports what was observed, organised by research question.
\Cref{ch:discussion} interprets those observations, relates them to the literature, draws out
the managerial implications and states the limitations. \Cref{ch:conclusions} answers the
research questions directly, summarises the contributions and identifies future work.
